\documentclass[8pt]{extarticle}
\usepackage[a4paper,margin=6mm,top=6mm,bottom=7mm]{geometry}
\usepackage{multicol,amsmath,amssymb}
\usepackage[T1]{fontenc}
\usepackage{xcolor}
\usepackage{enumitem}
\usepackage{array}
\setlength{\columnsep}{3.5mm}
\setlength{\columnseprule}{0.2pt}
\setlength{\parindent}{0pt}
\setlength{\parskip}{0.6pt}
\pagestyle{empty}
\setlist{nosep,leftmargin=3mm,label=\textbullet}
\newcommand{\sect}[1]{\par\vspace{1.5pt}\colorbox{black!12}{\parbox{\dimexpr\columnwidth-2\fboxsep}{\textbf{#1}}}\par\vspace{0.8pt}}
\newcommand{\sub}[1]{\textbf{#1}\;}
\renewcommand{\arraystretch}{1.05}
\begin{document}
\fontsize{7.9}{9.0}\selectfont\raggedright\setlength{\emergencystretch}{1em}
%% ==== SISI 1 ====
\begin{multicols*}{3}
\textbf{__JUDUL__} \hfill \textit{sisi 1 / 2, salin tulisan tangan}\par
\sect{0. Aturan main hitungan}
\begin{itemize}
\item Pembulatan, pembagi ($n-1$ atau $N$), kelas positif, konstanta ($\log_2x=\ln x/0{,}6931$), larangan dan cek silang khas mata kuliah ini.
\end{itemize}
\sect{a. Topik pertama}
\begin{itemize}
\item Definisi singkat, rumus dengan simbol slide, aturan keputusan, angka kontrol.
\end{itemize}
%% Tambah \sect{...} per topik. Utamakan yang mahal diturunkan ulang saat ujian.
\end{multicols*}
\newpage
%% ==== SISI 2 ====
\begin{multicols*}{3}
\textbf{__JUDUL__} \hfill \textit{sisi 2 / 2}\par
\sect{Urutan isian soal berantai}
\begin{itemize}
\item Urutan langkah tiap tipe soal worksheet, dari input sampai jawaban akhir.
\end{itemize}
\end{multicols*}
\end{document}
